\section{Instrument validity: what the persistence run inherits}
\label{sec:validity}

A retention measure compares two elicitations of the same item. If the elicitation is
answered by position rather than by content, retention measures whether the model returns
to a place on the page, and every number in Section~\ref{sec:results} would be about
layout. So the instrument is established first, and the two controls it forces are carried
into every episode.

\paragraph{Counterbalancing conceals the failure rather than fixing it.} This is the part
that cannot be worked around by averaging, and it is why the check has to be per item.
When a model answers purely by slot the two presentation orders return
$P(\text{option a}) = 0$ and $1$; their average is exactly $0.5$ and the order gap is
exactly $1.0$, so perfect balance and total position dependence are the same number. A
filter that inspects only the order-averaged split therefore selects \emph{for} the
artefact, and Figure~\ref{fig:gap-vs-mean} shows it doing so on real data: under
suppressed reasoning the pairs trace a sharp inverted V with its apex at $P = 0.5$, so the
closer a pair looks to balanced, the more completely its answer is determined by position.
A naive $0.3$--$0.7$ split filter accepts \pbNaiveOff{} pairs, of which \pbTrapOff{} have
an order gap above $0.5$. The order-averaged split is not a sufficient statistic; in the
failure case the two quantities are deterministically opposed.

\begin{figure}[t]
  \centering
  \includegraphics[width=0.72\textwidth]{gap_vs_mean.png}
  \caption{Order-averaged $P(\text{option a})$ against order gap on \pairModel{}, one
  point per pair per setting. The shaded column is what a naive $0.3$--$0.7$ split filter
  accepts; the darker block is the part of it that is position-driven. Perfect balance and
  total position dependence coincide at the apex.}
  \label{fig:gap-vs-mean}
\end{figure}

\paragraph{The two controls the persistence run inherits.} \textbf{Per-call reasoning is
enabled on every call.} With it suppressed, the mean order gap on \pairModel{} is
\pbOffMean{} (95\% CI [\pbOffLo{}, \pbOffHi{}]); enabling it moves the same pairs by
\pbDiff{} (95\% CI [\pbDiffLo{}, \pbDiffHi{}], \pbBoot{} resamples over pairs). The runner
asserts the setting and refuses to start otherwise. \textbf{Presentation order is fixed
within an episode}, so the re-elicitation shows the two options in the order the initial
elicitation used and a change of answer cannot be a change of layout; order is
counterbalanced across episodes and balanced within every pair $\times$ arm cell. Because
the answer token's log-probability saturates once reasoning is on, retention is scored on
the discrete choice rather than on log-probabilities.

\paragraph{The magnitude is model-specific, and one model refuses.} The same instrument on
the same \pbN{} pairs gives a reasoning-suppressed mean order gap of \pbOffMean{} on
\pairModel{} and \extGtwoOff{} on \texttt{gemini-2.5-flash} --- medians \extDsOffMed{} and
\extGtwoOffMed{}, so on the first model the typical pair is decided entirely by position
and on the second it is not affected at all --- while \texttt{gemini-3.5-flash} declines to
assign itself preferences in \extGthreeRefusalOn{}\% of trials and cannot be asked the
question at all. Nothing about the three models announced which would be which, so the
per-item order gap has to be measured on the model in use rather than inherited from a
published figure or a sister model. That analysis, the label-scheme controls that rule out
a quirk of the letter ``A'', and what both failure modes imply for welfare methods built on
elicited preferences are in Appendices~\ref{app:posbias}--\ref{app:welfare}.
